\documentclass{jpc}
\usepackage{booktabs,graphicx}
\usepackage{xr-hyper}
\makeatletter\g@addto@macro\UrlBreaks{\do\_\do\-}\makeatother
\externaldocument[M-]{jpc_main}
\renewcommand{\thesection}{S\arabic{section}}
\renewcommand{\theequation}{S\arabic{equation}}
\renewcommand{\thefigure}{S\arabic{figure}}
\renewcommand{\thetable}{S\arabic{table}}
\newcommand{\F}{\mathcal F}
\newcommand{\E}{\mathbb E}
\newcommand{\Prob}{\mathbb P}
\setlength{\emergencystretch}{3em}
\keywords{disclosure risk, sequential testing, supplementary material}
\begin{document}
\title[Supplement: adaptive certification of disclosure risk]{Supplementary material: Adaptive certification of attack-conditional disclosure risk across release programmes}
\author{Ronald Rethabile Malamule Mahomane}
\address{Department of Computing, Mathematical and Statistical Sciences, University of Namibia, Windhoek; School of Pathology, Faculty of Health Sciences, University of the Witwatersrand, Johannesburg; Basis Energy, Victoria, Australia}
\email{ron@basisenergygroup.com}
\begin{abstract}
This supplement gives the filtration and conditional validity argument, the allocation of the error budget, methods, reproduction instructions and the overlap disclosure for the main paper, together with the separate public-file search experiment.
\end{abstract}
\maketitle

\noindent References to theorems, propositions, assumptions, sections and tables without the prefix S are to the main paper.

\section{Filtration and conditional validity}\label{sec:s-filtration}
The global filtration $(\F_g)_{g\ge0}$ is generated by everything recorded up to global step $g$: development results, proposals, model updates, registrations, target draws and audit outcomes on every coordinate. Coordinate $i$ is registered at the stopping time $\sigma_i$, and its $b$th outcome is recorded at $\gamma_{i,b}$, with $\F_{\gamma_{i,b}-}=\F_{\gamma_{i,b}-1}$ the $\sigma$-field just before that draw. Between its own draws the e-process of coordinate $i$ is constant. At a draw, Assumption~\ref{M-ass:law} requires the conditional success probability under the null to be at least $\tau_i$ given everything recorded so far, including outcomes on other coordinates. With the identity in the proof of Theorem~\ref{M-thm:gate}, this makes the embedded process a nonnegative supermartingale in $(\F_g)$, that is, a test supermartingale \cite{shafer2011,ramdas2023}. This condition suffices for arbitrary global monitoring; it is not necessary. Independence of trials within a coordinate is not enough if correlated results on another coordinate reveal a target's hidden type and so change the conditional success probability.

\paragraph{Relaxed denominator.} Suppose an externally justified deviation $\kappa_i\ge0$ satisfies $\E[X_{i,b}\mid\F_{\gamma_{i,b}-}]\ge\tau_i-\kappa_i$ under the null. Then replace $\tau_i$ in the denominators of~\eqref{M-eq:bet} by $\tau_i-\kappa_i$, keep the policy tolerance $\tau_i$ in the null, and restrict the alternative to $q_{i,b}\in[0,\tau_i-\kappa_i)$. An alternative in $[\tau_i-\kappa_i,\tau_i)$ would make the expected factor exceed one for success probabilities just above $\tau_i-\kappa_i$. The change costs power. A drift statistic on features or a calibration metric does not by itself establish $\kappa_i$. Under cluster dependence, design trials around independent clusters or use an e-process valid for the cluster law; a reused target within a cluster is not a fresh audit unit.

\paragraph{Why marginal validity is not enough.} Let $H\sim\mathrm{Bernoulli}(1/2)$ be a hidden bit, drawn at registration and not observed, and set every attack result $X_b=H$ (the same target attacked repeatedly). Marginally $\E X_b=1/2$, but after one failure the conditional probability of a later success is zero. An auditor who treats each repeat as a fresh $\mathrm{Bernoulli}(1/2)$ trial can drive the evidence as high as it likes by attacking the same failed target again. Assumption~\ref{M-ass:law} fails, and no time-uniform error bound follows. Whether the assumption holds is a question of audit design, not of model calibration.

\paragraph{Weak and strict inequalities.} The null $r_i\ge\tau_i$ includes the boundary, so a coordinate whose risk equals its tolerance is certified with probability at most its error share. An agency that wants equality to count as safe can change the convention, but should not switch between strict and weak inequalities.

\paragraph{Proof details.} On $H_i$, each factor has conditional expectation at most one given $\F_{\gamma_{i,b}-}$, by conditions (b), (c) and Assumption~\ref{M-ass:law}. Because $(\theta_i,a_i)$, $\tau_i$ and $w_i$ are $\F_{\sigma_i}$-measurable and $r_i$ is a function of the registered scenario, $\mathbf 1_{H_i}$ is $\F_{\sigma_i}$-measurable; we set $w_i=0$ on $\{\sigma_i=\infty\}$. The conditional form of Ville's maximal inequality for nonnegative supermartingales \cite{ville1939} gives $\Prob(A_i\mid\F_{\sigma_i})\mathbf 1_{H_i}\le\alpha w_i\mathbf 1_{H_i}$. Summing unconditional bounds over $i=1,\ldots,M$ and letting $M\to\infty$ by monotone convergence, the pathwise bound $\sum_iw_i\le1$ gives the familywise bound. The argument holds when the rule for registering coordinate $i$ depends on outcomes of coordinates $1,\ldots,i-1$ and when audits are interleaved. It fails if $q_{i,b}$ is chosen after $X_{i,b}$ is seen, whatever it is later called.

\section{Allocation and zero-success calculation}\label{sec:s-allocation}
The allocation $w_i=6/(\pi^2i^2)$ is fixed at the $i$th registration and sums to one over countably many registrations; it is an instance of online Bonferroni \cite{tian2021}, the simplest case of alpha-spending \cite{foster2008}. An allocation that depends on the history at registration is also valid, provided the remaining budget is never overspent. The local threshold is $\log E_{i,b}\ge-\log(\alpha w_i)$, and each failure under a zero alternative adds $-\log(1-\tau_i)$. The implementation returns the trial at which the running sum first crosses the threshold, so it agrees with the gate in floating point: 29 trials for $\tau=0.12$, $\alpha=0.05$ and $w=1/2$, and 40 for $w=1/8$, where the threshold is $\log160$.

\paragraph{Cost weights and integer effects.} With planning drifts $\mu_i=\mu(p_i,q_i)$ from~\eqref{M-eq:drift} and trial costs $c_i$ for a frozen family, the expected-cost bound of Proposition~\ref{M-prop:weights} has the unique minimiser $w_i\propto c_i/\mu_i$; the zero-success case $p_i=q_i=0$ gives $\mu_i=u_i=-\log(1-\tau_i)$, and $u_i$ does not affect the minimiser. The rule gives more of the budget to costly audits and to coordinates planned close to their tolerance. It does not account for release utility, subgroup priority or scenarios added later, and a planning value that is too optimistic makes the audit longer, not invalid. Integer ceilings make the true objective a step function; neighbouring allocations can be enumerated after computing the continuous one, but must be fixed before outcomes are seen.

\paragraph{Reserve for later coordinates.} Choose $0<\beta<1$. Give $\beta$ to the known family by the cost rule and reserve $(1-\beta)\,6/(\pi^2j^2)$ for the $j$th coordinate introduced later. The sequence is summable on every path, even when a later coordinate is triggered by earlier outcomes. A reserve is not protection against an unregistered harm.

\section{Numerical risk bounds and the audited margin}\label{sec:s-bounds}
For a fixed sample size $B$ and $Y$ independent Bernoulli successes, the exact one-sided upper bound $\overline{R}$ solves $\sum_{m=0}^{Y}\binom Bm \overline{R}^m(1-\overline{R})^{B-m}=\lambda$; at $Y=0$, $\overline{R}=1-\lambda^{1/B}$ \cite{clopper1934}. For a fixed family, use $\lambda=\alpha/M$ and a union bound \cite{dunn1961}. If one inspects repeatedly and stops when the fixed-$B$ bound is favourable, that coverage claim no longer holds. A time-uniform upper confidence sequence, obtained by inverting a family of valid sequential tests indexed by hypothetical risks, can supply $\overline{R}_i$ for the audited margin \cite{darling1967,lai1976,howard2020,howard2021,waudby2024} $z_i=1-\overline{R}_i/\tau_i$. The e-process at a single $\tau_i$ tests eligibility at that tolerance and does not output such a sequence.

A report should say, for every figure, whether it is a point estimate, a model prediction, an upper confidence bound or a tolerance. The crossing indicator of the gate belongs in a separate field with its level $\alpha w_i$, the number of draws and the attack protocol. Reporting an e-process as a risk estimate is an error.

\section{A proposal model for bets}\label{sec:network}
Any predictor may supply the bets $q_{i,b}$, provided each is fixed before $X_{i,b}$ is observed and clipped to $[0,\tau_i)$; validity does not depend on its accuracy (Section~\ref{M-sec:bets}). The repository contains one such predictor, SENTRY-R, which fits a logistic regression or a small ReLU network to counted development attacks by the binomial log likelihood and ranks candidate releases by their largest predicted risk-to-tolerance ratio. It keeps three streams apart: scenarios proposed by a language model or a reviewer, which can enlarge the threat model but supply no outcomes; development attacks, which train the predictor; and final audit outcomes, which alone can certify. A risk probability asserted by a language model is not a Bernoulli outcome. On a synthetic programme law, over 10 seeds with both models trained to convergence, logistic regression had the lower test log loss in all 10 (0.2556 against 0.2576 for a three-layer network). This benchmark concerns SENTRY-R. The network of Section~\ref{M-sec:neural-run} is the neural release pipeline, \nolinkurl{privacy_risk.neural_pipeline}; its development counts, fitted weights, frozen plans and final outcomes are in \nolinkurl{artifacts-neural/neural_pipeline.json}. Training equations and the benchmark are in \nolinkurl{docs/SENTRY_R.md}.

\section{Empirical methods and provenance}\label{sec:s-methods}
\paragraph{Carried-over results.} The synthetic population of \cite{mahomane2026} has 200,000 units and a five per cent release sample; its four-release experiment draws 20,000 targets and intersects the candidate sets of four 3-anonymous key projections. The synthetic credit register has six monthly vintages; for each policy the linkage evaluation uses 4,000 attacks on independently drawn targets, with one exposure reading per vintage corrupted by Gaussian noise. Three harms are evaluated beyond record identification: precision of an alleged lender relationship, disclosure of arrears status, and differencing. The main paper uses identification only. The public Census sample holds 29,501 records; neither it nor the other public file used in \cite{mahomane2026} supplies the population from which it was sampled. The release and pooled columns of Table~\ref{M-tab:carried} and the grid used to choose the minimum cell count come from separate target draws.

\paragraph{Toy audit.} \nolinkurl{scripts/adaptive_audit_comparison.py} uses seed 2709, 2,000 replications per law, at most 300 draws, $\alpha=0.05$, $w=1/8$, $\tau=0.12$, $q=0.03$ and true risks $0.04$ and $0.12$. The first-passage event is $\max_{b\le300}\log E_b\ge\log160$. The fixed comparator uses only the count at draw 300 and its exact one-sided bound at level $\alpha w$. In the adaptive run a new stream with risk $0.12$ is created only if an independent stream with risk $0.04$ has crossed. Monte Carlo standard errors in Table~\ref{M-tab:new-audit} are $\sqrt{\hat p(1-\hat p)/2000}$.

\paragraph{Register.} The synthetic register of \nolinkurl{src/privacy_risk/credit.py} draws 20,000 borrowers, 60 lenders and 32,581 loans over six months. A share of loans is originated after the first month, each loan leaves with a monthly hazard of 0.04 and does not return, and the stock stays near 27,000. This replaces the panel of \cite{mahomane2026}, in which every loan was forced to be live in the first month and a third of loans left in the second month and came back in the third, which gave the history adversary an artificial signal. Exposure is lognormal within size class and constant across vintages.

\paragraph{Credit-register audit.} \nolinkurl{scripts/adaptive_audit_chain.py} runs the release chain over the floors 300 to 1 with the three declared attacks, the bet $\min\{(Y+1/2)/(B+1),0.11\}$, at most 3,000 fresh targets per coordinate and threshold $1/\alpha$, using \nolinkurl{AdaptiveAudit.register_release_chain}. The five other designs of Section~\ref{M-sec:credit} are applied to the same streams. True risks come from separate draws of 100,000 targets, on seeds disjoint from the audit seeds. The operating-characteristics check runs the release chain 400 times on Bernoulli streams at the estimated risks, and at risks raised by 0.01 and 0.02.

\paragraph{Factor-by-factor floors.} \nolinkurl{scripts/programme_robustness.py} chooses the floor for ten adversaries (the five of Table~\ref{M-tab:robust}, the insider with log-exposure error 0.01, and four with exposure published to one significant figure) by fixed-sequence testing on one draw of 4,000 targets per level and confirms it on a second draw; where confirmation fails, the next stricter level that passes is used, which keeps the error bound because every stricter level also passed at selection. It audits the largest size class on its own stratified draw, counts the loans whose key and published exposure are unique (the relink share under rotating pseudonyms), records the share of loans, exposure, large-firm loans and multi-lender borrowers kept at each floor, gives floors at tolerances of 0.05, 0.09, 0.12 and 0.20, and repeats the choice over 20 register draws and reading standard deviations of 0.3, 0.6 and 1.0 with 2,000 targets per level. At the declared noise of 0.6 the floors over register draws were 1 for one release, 40 to 75 (median 75) for pooling with the history, and 100 to 150 (median 150) for the full attack; at 0.3 they were 20, 150 and 300, with no level passing in 3 of 20 draws; at 1.0 they were 1, 40 and 75. \nolinkurl{scripts/robustness_table.py} writes the tables.

\paragraph{Audit-process attacks.} \nolinkurl{scripts/audit_process_attacks.py} simulates 20,000 Bernoulli streams per risk with one coordinate at weight one, and 2,000 runs, each of at most 300 draws, of an operator that bets just below the tolerance after each success and zero after each failure. The rates for ten and a hundred tries are computed as $1-(1-f)^{\rm tries}$ from the single-try rate $f$, not simulated. The controls it models are in \nolinkurl{src/privacy_risk/adaptive_audit.py}: \texttt{commit} and \texttt{target\_seed} derive targets from a committed secret and a beacon round, \texttt{CommittedCoordinate} records each bet in an \texttt{AuditLedger} before the attack runs, and the ledger is hash-chained. The script simulates the effect of these controls on Bernoulli streams; it does not call them.

\section{Algorithm and reproduction}\label{sec:s-algorithm}
\begin{enumerate}
\item Fix the harm taxonomy and the pooled attack protocol, and keep the release history and data versions.
\item Register each coordinate with its scenario, tolerance and weight $w_i$, drawn from the remaining budget. Registered fields cannot be changed.
\item Before each attack choose $q_{i,b}\in[0,\tau_i)$ from earlier information only. Record the target draw, outcome, bet, sampling-unit identifier and attack code digest.
\item Add $\log(q/\tau)$ to the log evidence on a success and $\log((1-q)/(1-\tau))$ on a failure; a success with $q=0$ sets it to $-\infty$. Each step adds the log likelihood ratio of Bernoulli$(q)$ to Bernoulli$(\tau)$, as in the sequential probability ratio test \cite{wald1945}.
\item Certify at the first crossing of $-\log(\alpha w_i)$ and accept no further outcomes. Call a release eligible only when every declared coordinate is certified, and record the owner and the expiry of the decision.
\end{enumerate}

The following commands, run from the repository root, install the repository, test the gate, recreate Tables~\ref{M-tab:new-audit}, \ref{M-tab:chain-audit}, \ref{M-tab:robust}, \ref{M-tab:tolerance-floors} and~\ref{M-tab:audit-process}, and rerun the pipeline of \cite{mahomane2026} that produces the carried-over results:
\begin{flushleft}\small
\texttt{python -m pip install -e '.[dev]'}\\
\texttt{python -m pytest tests/test\_adaptive\_audit.py -q -o addopts=''}\\
\texttt{python scripts/adaptive\_audit\_comparison.py}\\
\quad\texttt{--output paper/publication\_versions/adaptive\_audit\_toy.json}\\
\texttt{python scripts/adaptive\_audit\_chain.py}\\
\texttt{python scripts/programme\_robustness.py}\\
\texttt{python scripts/robustness\_table.py}\\
\texttt{python scripts/audit\_process\_attacks.py}\\
\texttt{python -m privacy\_risk.pipeline --config configs/paper.json}\\
\quad\texttt{--output artifacts-paper}
\end{flushleft}
The code checks inputs and the error budget; it cannot check that attack draws satisfy Assumption~\ref{M-ass:law} or that the declared coordinate set is complete.

\section{Relation to earlier work by the author}\label{sec:s-overlap}
Table~\ref{M-tab:carried} and the four-release and public-file figures in Section~\ref{M-sec:carried} of the main paper were first reported in \cite{mahomane2026}, an author manuscript whose earlier version was presented at the Banco de Portugal and IFC Workshop on Data Security and Privacy, Lisbon, September 2026. That study did not run the gate. Theorem~\ref{M-thm:gate}, Propositions~\ref{M-prop:cost} to~\ref{M-prop:release}, the toy audit, the credit-register audit on the revised register, the factor-by-factor floors and the audit-process measurements are new to the present article.

\section{Multiscale proposal models and audited distribution shift}
\label{sec:multiscale-extension}
This extension separates three mathematical objects: the unknown attack law,
the fitted proposal function, and a confidence bound from new attack trials.
Its results do not replace the registered sequential gate. In particular,
a derivative of a neural network describes its response to an input change;
it does not bound an adversary's true success probability. The new software is
\texttt{privacy\_risk.multiscale}. All new numerical demonstrations use synthetic
laws and do not modify the earlier public-data results.

\subsection{An arbitrary dimensional risk vector}
Let $x\in\mathcal X=\prod_{j=1}^d[\ell_j,u_j]$ describe a release and its
registered threat model, with $u_j>\ell_j$. For coordinate $i\in\{1,\ldots,m\}$,
let $A_i(Z,x)\in\{0,1\}$ be the declared attack success event. Define
\begin{equation}
 r_i(x;P_i)=\int A_i(z,x)\,dP_i(z),\qquad
 z_i(x)=1-r_i(x;P_i)/\tau_i,\quad 0<\tau_i<1.
\end{equation}
A coordinate includes its target population, information available to the
attacker, attack algorithm and harm. These components cannot change silently
when the feature vector changes. Eligibility is the intersection of all
registered inequalities $r_i\le\tau_i$, rather than an average over harms.
The scalar $\min_i z_i$ records the limiting margin, but a report retains the
whole vector so that a favourable component cannot conceal another violation.

\subsection{Finite functional wavelet dictionary}
Normalize $v_j=(x_j-\ell_j)/(u_j-\ell_j)$ and define
\begin{equation}
 \psi(t)=(1-t^2)e^{-t^2/2},\qquad
 \psi'(t)=(t^3-3t)e^{-t^2/2}.
\end{equation}
For levels $a=0,\ldots,J-1$ and centres
$c_{ak}=(k+1/2)2^{-a}$, $k=0,\ldots,2^a-1$, the feature map is
\begin{equation}
 \Phi_J(x)=\big(v_j,\psi(2^a(v_j-c_{ak}))\big)_{j,a,k}
 \in\mathbb R^{D},\qquad D=d\,2^J.
\end{equation}
The dictionary captures localized changes at several scales and retains the
normalized identity coordinates. It is neither an orthogonal wavelet transform
nor an exact interpolant. On a bounded interval the atoms need not be orthogonal.
An exact interpolation claim would require specified nodes, an invertible
collocation matrix and a conditioning estimate, none of which follows from
using a wavelet-shaped activation. Tensor products of all atoms cost
exponentially in $d$; the coordinatewise dictionary avoids that construction.
Cross-feature interactions arise in the dense network. Its computational benefit
is a choice of representation, not a universal approximation-error guarantee.

Put $s_j=(u_j-\ell_j)^{-1}$. Each feature depends on a single raw input, so the
feature Jacobian has orthogonal columns. Since
$|\psi'|\le C_\psi=3\sqrt3 e^{-3/2}+3e^{-1/2}$,
\begin{equation}
 \|D\Phi_J(x)\|_2\le L_\Phi=
 \max_j s_j\left(1+C_\psi^2\sum_{a=0}^{J-1}2^a4^a\right)^{1/2}.
\end{equation}
The bound follows by separately maximizing $|t|^3e^{-t^2/2}$ and
$3|t|e^{-t^2/2}$. Integration on a line segment in the convex box establishes
$\|\Phi_J(x)-\Phi_J(y)\|_2\le L_\Phi\|x-y\|_2$.
Increasing resolution increases $L_\Phi$ and may worsen robustness; fine scales
are not automatically an improvement.

\subsection{Deep counted likelihood and explicit sensitivities}
For column-vector notation let $h_0=\Phi_J(x)$,
$h_l=\operatorname{ReLU}(W_lh_{l-1}+b_l)$, $l=1,\ldots,L$, and
$f_\theta(x)=\sigma(W_{L+1}h_L+b_{L+1})\in(0,1)^m$.
For development counts $k_{ni}$ out of $b_{ni}$ trials, fit
\begin{equation}
 \mathcal L(\theta)= -\frac{1}{\sum_{n,i}b_{ni}}
 \sum_{n,i}\{k_{ni}\log f_{ni}+(b_{ni}-k_{ni})\log(1-f_{ni})\}
 +\frac\lambda2\sum_l\|W_l\|_F^2.
\end{equation}
Binomial counts retain different trial denominators. Treating all observed
fractions as equally precise generally changes this objective. If the counts
are dependent, the expression is a working loss, not an independent binomial
likelihood. Training, early-stopping validation and final audits must use
separate groups; neither LLM proposals nor final certification trials provide
training labels.

Away from ReLU kinks the raw-input derivative is
\begin{equation}
 Df_\theta(x)=D_\sigma W_{L+1}D_LW_L\cdots D_1W_1D\Phi_J(x),
\end{equation}
where $D_l$ contains activation indicators and
$D_\sigma=\operatorname{diag}(f_i(1-f_i))$. The implementation stores transposed
weights because observations are rows. At ReLU kinks it uses the zero
subgradient. For logits clipped to $[-35,35]$, it additionally sets the
logit derivative to zero outside that interval. Thus it differentiates the
actual deployed predictor, rather than an unclipped surrogate.

Because ReLU and clipping are one-Lipschitz and $\sigma'\le1/4$,
\begin{equation}
 \|f_\theta(x)-f_\theta(y)\|_2\le
 L_f\|x-y\|_2,\quad L_f=\tfrac14 L_\Phi\prod_{l=1}^{L+1}\|W_l\|_2.
\end{equation}
This global bound is conservative; local Jacobians can be much smaller.
For coordinate $i$, the proposal margin derivative is
$D\widehat z_i=-Df_i/\tau_i$. A useful sensitivity report gives raw units,
feature ranges, point derivatives and the spectral bound. It never labels
$\widehat z_i$ an audited margin. If a separate uniform approximation guarantee
$\sup_x|r_i(x)-f_i(x)|\le\epsilon_i$ were established, the triangle inequality
would imply $r_i(y)\le f_i(x)+L_f\|y-x\|+\epsilon_i$. A held-out mean log loss
does not establish that uniform guarantee.

\subsection{A finite-family audit with a stated shift radius}
Fix the family, sample sizes and weights before final outcomes are observed.
For coordinate $i$, let $K_i\sim\operatorname{Bin}(n_i,p_i)$ under audit law
$P_i$, and let $Q_i$ be the intended deployment attack law. Assume an externally
justified total-variation bound
\begin{equation}
 \operatorname{TV}(P_i,Q_i)=\sup_E|P_i(E)-Q_i(E)|\le\rho_i.
\end{equation}
For weights $w_i>0$, $\sum_iw_i\le1$, use the exact one-sided
Clopper--Pearson bound
\begin{equation}
 U_i=\begin{cases}
 \operatorname{Beta}^{-1}(1-\alpha w_i;K_i+1,n_i-K_i),&K_i<n_i,\\
 1,&K_i=n_i.
 \end{cases}
 \quad U_i^Q=\min\{1,U_i+\rho_i\}.
\end{equation}
With probability at least $1-\alpha$, all deployment risks obey
$r_i(Q_i)\le U_i^Q$. To prove this, each binomial bound fails with probability
at most $\alpha w_i$; the union bound controls their simultaneous failure.
On the complementary event, the definition of total variation applied to the
attack-success event gives $r_i(Q_i)\le p_i+\rho_i\le U_i+\rho_i$.
No independence across coordinates is needed. Independent Bernoulli trials
within each coordinate, a fixed sample size and a fixed target law are needed.
Random data-dependent radii require their own simultaneous error budget.
A feature displacement and a model Lipschitz bound do not justify a TV radius.

The shifted audited position is $z_i^Q=1-U_i^Q/\tau_i$; all must be nonnegative.
This is a conditional statistical guarantee given the external shift model,
not proof that the radius is true or that unregistered harms are absent.
At zero successes, $U_i=1-(\alpha w_i)^{1/n_i}$, so passage requires
\begin{equation}
 n_i\ge\left\lceil\frac{\log(\alpha w_i)}
 {\log(1-\tau_i+\rho_i)}\right\rceil,\qquad 0\le\rho_i<\tau_i.
\end{equation}
If $\rho_i\ge\tau_i$, zero observed successes cannot supply a finite passing
budget. Repeatedly inspecting these fixed-size bounds and stopping at passage
invalidates their stated coverage; use the existing sequential gate instead.

\subsection{Audit cost and scaling}
For positive costs $c_i$, set $a_i=c_i/[-\log(1-\tau_i+\rho_i)]$.
Ignoring integer ceilings, minimize
$C(w)=\sum_i a_i\log(1/(\alpha w_i))$ over positive weights summing to one.
The Lagrangian derivative is $-a_i/w_i+\nu=0$; summing gives
\begin{equation}
 w_i^*=a_i/\sum_ja_j.
\end{equation}
The Hessian is diagonal with entries $a_i/w_i^2>0$, proving a unique minimum.
The integer plan rounds upward and checks the exact zero-success inequality.
It is a planning rule conditional on zero successes, not a promised audit cost
under positive risk and not an optimizer of the discontinuous integer objective.
Freeze weights from costs and declared tolerances before audit outcomes.

Feature construction costs $O(Nd2^J)$. A dense network forward pass costs
$O(N\sum_l d_{l-1}d_l)$; storing every raw-input sensitivity additionally costs
$O(Nmd)$ for the output alone. The implementation supports arbitrary $d$ but
limits $J$ to eight to avoid accidental dictionary growth. Dimension-free
statistical safety does not imply dimension-free computational or predictive
accuracy. LLM scenario injections continue through the existing review and
measured-development-evidence pipeline; modifying a dictionary or threat law
requires a new proposal-model version and new final audits.


\def\publicsearchtable{../generated/table_public_search}
\def\publicsearchllmtable{../generated/table_public_search_llm}
\section{Separate public-file search experiment}\label{sec:public-search}

This additional exploratory experiment measures the search for high-risk
scenarios, not the safety certification of a release. It does not alter the
synthetic-register floors or the gate results reported in the main text. A
separate working draft, \emph{Measured Scenario Injection for Privacy Audits:
Search Coverage, Evidence Integrity, and Independent Confirmation}, develops
the experiment and its relation to an LLM-assisted proposal pipeline. A
fifth arm, added after the four baseline arms were run and reported below,
replays proposals from a small open-weight language model.

The source is the 29,501-row public \texttt{census2000} extract distributed
with \texttt{wooldridge} 0.5.0. We select the first 30 encoded state groups
containing at least 100 rows, before observing attack outcomes. An attacker
uses a matching-key prior and an independently perturbed log-income reading
to identify a uniformly sampled record within each released empirical file.
Targets are drawn with replacement, and the auxiliary noise is fresh. The
declared choices vary PUMA coarsening, education coarsening, minimum released
cell size and auxiliary-reading precision. Settings retaining fewer than 30
records are excluded. These measurements are conditional on already public
records and a simulated auxiliary-reading law; they do not identify risk in
an undisclosed population.

Each state shares 16 warm-start settings and eight early-stopping settings
in a separate state, each measured with 64 trials. Each of random, informed
heuristic, scrambled Sobol and neural-guided search then uses 32 additional
settings with 64 trials each. Every selected setting receives 300 fresh
confirmation trials. The total is 1,443,840 attack trials. Neural search
fits a width-16, two-hidden-layer model four times per state; the other arms
do not fit a model, so evaluator cost is matched but total compute is not.

A selected setting is confirmed high risk only if a one-sided exact binomial
lower limit exceeds 0.2. Its per-test level is $0.05/128$, accounting for all
32 selected settings in each of four arms within one state. Conditional on
the complete search history, a union bound controls the probability of any
false high-risk confirmation at 0.05 within that state, under the stated
fresh-draw law. This is not simultaneous control across all 30 states, and
lower-bound confirmation of excessive risk is not upper-bound certification
of safety.

The primary endpoint is the number of distinct confirmed cells in a fixed
24-cell taxonomy: three PUMA levels, education width above or below two,
minimum cell count equal to one or larger, and auxiliary standard deviation
above or below 0.5. The primary contrast is neural search minus Sobol;
intervals resample paired state blocks 10,000 times. They are descriptive
intervals for this convenience subset, not national design-based inference.

\input{\publicsearchtable}

At the initial 600-epoch cap, neural search covers 3.77 fewer confirmed cells
than Sobol per state (paired interval $[-4.43,-3.07]$), but confirms more
individual high-risk settings: 31.57 versus 22.03 of 32. The heuristic
confirms all 32 on average while covering fewer cells than Sobol. Four of
the 120 neural fits have their best epoch at the cap. A sensitivity run with
a 2,000-epoch cap retains the neural mean of 10.07 cells and a difference of
$-3.77$ (interval $[-4.40,-3.07]$). It reuses the same data and non-neural
draws; it is not independent replication or a convergence proof.

\paragraph{A language model as the fifth arm.} The saved per-state prompt
was then sent, byte for byte, to \texttt{qwen3:4b}, a 4.0B-parameter
open-weight model (Q4\_K\_M quantisation) run locally through Ollama
0.40.9 on the author's machine, so no account, key or charge was
involved and nothing left the machine. Thinking was off and decoding was
constrained to the JSON shape the prompt asks for; the tuples the model chose
were still checked afterwards for count, distinctness and catalogue membership.
Each state was allowed three attempts with distinct seeds. The model returned a
valid list of 32 tuples for 20 of 30 states
(7 at the first attempt); 50 of
70 attempts were rejected, 43
for repeated or off-catalogue tuples and the rest for repeated identifiers. All
five arms were then rerun on exactly those 20 states
(1,195,520 attack trials), with the per-test level at $0.05/160$.
Table \ref{tab:public-search-llm} gives the result. The model confirmed
27.60 of 32 settings per state, between the two
yield-maximising arms and the two space-filling arms, but covered only
6.05 cells against 13.65 for Sobol and
10.10 for neural search; the paired difference from Sobol is
-7.60 cells (interval $[-8.30,-6.80]$).
The reason is visible in the proposals: in 19 of
20 states every tuple the model chose sits at the finest
geography, and the lists follow the catalogue's own order. The model copied the
front of the list; it did not search. This is one small model under one fixed
prompt. It settles that such a model adds no coverage here; it does not settle
what a larger model, a prompt that names the coverage goal, or a worked example
would do, and any such change must be declared before it is run.

\input{\publicsearchllmtable}

These results support distinguishing targeting yield from coverage, not
the superiority of either neural models or LLMs. The evaluator executes only
one attack family. Prose proposals cannot create new executable attacks
without a separately specified evaluator. Reproduction commands are
\texttt{python scripts/public\_injector\_search.py --states 30} and
\texttt{python scripts/public\_injector\_search.py --states 30 --epochs 2000 --output artifacts-public-search-2000};
the language-model arm is \texttt{python scripts/collect\_llm\_pilot\_local.py --states <codes>}
followed by \texttt{python scripts/public\_injector\_search.py --states 30 --states-from-manifest --llm-replays artifacts-llm-local-30/manifest.json --output artifacts-public-search-llm-local-30}.
The source-data hash, state-level measurements and design are in
\texttt{artifacts-public-search/} and
\texttt{artifacts-public-search-2000/}; the model's proposals, every attempt and the
five-arm rerun are in \texttt{artifacts-llm-local-30/} and
\texttt{artifacts-public-search-llm-local-30/}.


\begin{thebibliography}{99}\small
\bibitem{clopper1934} Clopper, C. J. and Pearson, E. S. (1934). The use of confidence or fiducial limits illustrated in the case of the binomial. \emph{Biometrika} 26(4), 404--413.
\bibitem{darling1967} Darling, D. A. and Robbins, H. (1967). Confidence sequences for mean, variance, and median. \emph{Proceedings of the National Academy of Sciences} 58(1), 66--68.
\bibitem{dunn1961} Dunn, O. J. (1961). Multiple comparisons among means. \emph{Journal of the American Statistical Association} 56(293), 52--64.
\bibitem{foster2008} Foster, D. P. and Stine, R. A. (2008). $\alpha$-investing: a procedure for sequential control of expected false discoveries. \emph{Journal of the Royal Statistical Society, Series B} 70(2), 429--444.
\bibitem{howard2020} Howard, S. R., Ramdas, A., McAuliffe, J. and Sekhon, J. (2020). Time-uniform Chernoff bounds via nonnegative supermartingales. \emph{Probability Surveys} 17. doi:10.1214/18-PS321.
\bibitem{howard2021} Howard, S. R., Ramdas, A., McAuliffe, J. and Sekhon, J. (2021). Time-uniform, nonparametric, nonasymptotic confidence sequences. \emph{Annals of Statistics} 49(2), 1055--1080.
\bibitem{lai1976} Lai, T. L. (1976). On confidence sequences. \emph{Annals of Statistics} 4(2). doi:10.1214/aos/1176343406.
\bibitem{mahomane2026} Mahomane, R. R. M. (2026). Governing privacy risk across repeated statistical releases: probability models, programme tolerances and accountable decisions. Author manuscript; an earlier version was presented at the Banco de Portugal and IFC Workshop on Data Security and Privacy, Lisbon, September 2026.
\bibitem{ramdas2023} Ramdas, A., Gr\"unwald, P., Vovk, V. and Shafer, G. (2023). Game-theoretic statistics and safe anytime-valid inference. \emph{Statistical Science} 38(4), 576--601.
\bibitem{shafer2011} Shafer, G., Shen, A., Vereshchagin, N. and Vovk, V. (2011). Test martingales, Bayes factors and $p$-values. \emph{Statistical Science} 26(1), 84--101.
\bibitem{tian2021} Tian, J. and Ramdas, A. (2021). Online control of the familywise error rate. \emph{Statistical Methods in Medical Research} 30(4), 976--993.
\bibitem{ville1939} Ville, J. (1939). \emph{\'Etude critique de la notion de collectif}. Gauthier-Villars, Paris.
\bibitem{wald1945} Wald, A. (1945). Sequential tests of statistical hypotheses. \emph{Annals of Mathematical Statistics} 16(2), 117--186.
\bibitem{waudby2024} Waudby-Smith, I. and Ramdas, A. (2024). Estimating means of bounded random variables by betting. \emph{Journal of the Royal Statistical Society, Series B} 86(1), 1--27.
\end{thebibliography}
\end{document}
